\documentclass[10pt,a4paper]{amsart}
\usepackage[margin=19mm]{geometry}
\usepackage{amsmath,amssymb,mathtools}
\usepackage[scr=rsfs,cal=euler]{mathalfa}
\usepackage{xfrac}
\usepackage[unicode,hidelinks,bookmarks=false]{hyperref}
\newcommand{\ie}{i.e.\ }
\newcommand{\cf}{cf.\ }
\newcommand{\resp}{respectively\ }
\newcommand{\Subsection}{\subsection}
\providecommand{\nobreakdash}{\nobreak\hbox{-}}
\setlength{\emergencystretch}{2em}
\multlinegap0pt
\usepackage{amssymb}
\usepackage{xcolor}
\newtheorem{theorem}{Theorem}
\title{Abelian-normal decimal expansions}
\author{John M. Campbell}
\date{Source: arXiv:2603.04396v2 (CC BY 4.0)}
\begin{document}
\maketitle

\noindent\emph{Source and licence.} Abelian-normal decimal expansions. Version arXiv:2603.04396v2 (CC BY 4.0), distributed by arXiv under the Creative Commons Attribution 4.0 International licence, http://creativecommons.org/licenses/by/4.0/, as recorded in the arXiv licence metadata for this exact version and shown on the abstract page https://arxiv.org/abs/2603.04396v2. Full licence notice: \texttt{licenses/arxiv-2603.04396-CC-BY-4.0.txt}.\par\smallskip

\noindent\textit{Editorial scope.} Counted unit: the article's theorem theoremXi10 (source0.tex lines 184--187). The article's complete original proof of it is delivered with it (source0.tex lines 189--389, one \texttt{proof} environment of 201 lines). No mathematics is written, repaired or re-derived: the statement and the proof are the article's own lines, in the article's own notation, and no retained line is substituted or re-flowed.  Retained uncounted as the source-local context the proof needs: lines 157--165.  Nothing was omitted from any retained span, so the package carries no omission record.  No retained line contains a citation, so the wrapper carries no bibliography and no bibliography entry.  No retained line contains a figure environment, an image-inclusion command or drawing code, so no image is delivered and the package declares no asset dependency.  Editorial formatting only, in addition: an AMS \texttt{amsart} preamble that restates the article's own declarations and macros used by the retained lines, namely the environments \texttt{theorem} on separate counters, together with the standard packages those lines need.  The retained environments print in this document as Theorem 1 in order of appearance, whereas the article prints the same items with the numbering of its own full document.  The complete native source of the article as distributed in the arXiv e-print for this exact version is bundled unchanged as \texttt{sources/195741/source0.tex}.
\begin{equation}\label{definephi}
 \phi(10a+b)=
 \begin{cases}
 01, & (a,b)=(1,0),\\
 12, & (a,b)=(2,1),\\
 20, & (a,b)=(0,2),\\
 ab, & \text{otherwise.}
 \end{cases}
\end{equation}

\begin{theorem}\label{theoremXi10}
 The constant $\Xi_{10}$ is abelian-normal in base $10$ but not normal
 in base $10$.
\end{theorem}

\begin{proof}
   Let $U$ be uniformly distributed over the labels in   $\{0,1,\ldots,99\}$, and write $XY=\phi(U)$,     with the understanding that $X$ and  
  $Y$, respectively, 
 are equal to the first and second digits (written as single-character words) of $\phi(U)$, 
 and where $XY$ denotes the concatenation of these words. 
 It follows in an immediate way from
\eqref{definephi} that
\begin{equation}\label{PrXYab}
 \operatorname{Pr}(XY=ab)=
 \begin{cases}
 \dfrac{2}{100}, & ab\in\{01,12,20\},\\
 0, & ab\in\{10,21,02\},\\
 \dfrac{1}{100}, & \text{otherwise.}
 \end{cases}
\end{equation}
 While the distribution in \eqref{PrXYab} 
 is not uniform, each of the associated 
 coordinate distributions is uniform, with 
\begin{equation}\label{twouniform}
 \operatorname{Pr}(X=d)=\operatorname{Pr}(Y=d)=\frac{1}{10} 
\end{equation}
 for $ d \in \{ 0, 1, \ldots, 9 \}$. 
 Indeed, in the first coordinate, the three exceptional substitutions 
 (indicated in the definition of $\phi$ in \eqref{definephi}) 
 replace $1$ by $0$, $2$ by $1$, and $0$ by $2$, whereas, in the second
coordinate, these substitutions replace $0$ by $1$, $1$ by $2$, and $2$ by $0$.
Consequently, the number of occurrences of each digit in either
coordinate is unchanged.

For commuting indeterminates $z_0,z_1,\ldots,z_9$, write
$\mathbf z=(z_0,z_1,\ldots,z_9)$ and set
\begin{equation*}
 S(\mathbf z)=\frac{z_0+z_1+\cdots+z_9}{10}.
\end{equation*}
We also set
\begin{equation*}
 P_j(\mathbf z)
 =\frac{1}{100}\sum_{u=0}^{99}z_{\phi_j(u)} 
\end{equation*}
 for $j \in \{ 1, 2 \}$
 and 
 $$ Q(\mathbf z)
 =\frac{1}{100}\sum_{u=0}^{99}
 z_{\phi_1(u)}z_{\phi_2(u)}.
 $$
 From \eqref{twouniform}, we find that 
\begin{equation}\label{singlecoordinateidentities}
 P_1(\mathbf z)=P_2(\mathbf z)=S(\mathbf z).
\end{equation}
Moreover, each of the three exceptional substitutions in
\eqref{definephi} reverses a pair of digits. Since the given 
indeterminates commute, the corresponding monomial is unchanged.
 As a consequence, we have that 
\begin{equation*}
 Q(\mathbf z)
 =\frac{1}{100}\sum_{a,b=0}^{9}z_az_b
 = S^2(\mathbf z).
\end{equation*}

 We next record a consequence of the base-$100$ normality of
$c_1c_2\cdots$ that is required for our purposes. For every fixed positive
integer $s$ and every function
$F\colon\{0,1,\ldots,99\}^{s}\longrightarrow\mathbb C$, we claim that 
\begin{equation}\label{normalaveraging}
 \lim_{N\to\infty}\frac{1}{N}
 \sum_{i=1}^{N}F(c_i,c_{i+1},\ldots,c_{i+s-1}) =\frac{1}{100^{s}}
 \sum_{u_1,\ldots,u_s=0}^{99}F(u_1,\ldots,u_s). 
\end{equation} 
 This may be seen by grouping the terms on the left-hand side 
 of \eqref{normalaveraging}
 according to each possible
length-$s$ block, as every such block is of limiting frequency $100^{-s}$. 
 
 Now, write $$  x_1 x_2 x_3 \cdots=\phi(c_1)\phi(c_2)\phi(c_3)\cdots, $$ so that
\begin{equation}\label{xcoordinates}
 x_{2i-1}=\phi_1(c_i)
 \quad\text{and}\quad
 x_{2i}=\phi_2(c_i) 
\end{equation} 
 for each positive integer $i$. 
For a finite decimal word $w=w_1w_2\cdots w_m$, let
\begin{equation}\label{defineM}
 M(w;\mathbf z)
 =\prod_{j=1}^{m}z_{w_j}
 =\prod_{d=0}^{9}z_d^{\nu_d(w)},
\end{equation}
where $\nu_d(w)$ is the number of occurrences of the digit $d$ in $w$.
All polynomial limits below are understood coefficientwise.

 If $m=2r$ (for $r \geq 1$), then \eqref{normalaveraging} gives 
\begin{equation}\label{oddstartevenlength}
 \lim_{N\to\infty}\frac{1}{N}
 \sum_{i=1}^{N}
 M(x_{2i-1}\cdots x_{2i+2r-2};\mathbf z) 
 = Q^r(\mathbf z) 
 = S^{2r}(\mathbf z).
\end{equation}
 If $m=2r+1$ (for $r \geq 0$), then \eqref{normalaveraging} gives 
\begin{equation}%keepitlabeled
 \lim_{N\to\infty}\frac{1}{N}
 \sum_{i=1}^{N}
 M(x_{2i-1}\cdots x_{2i+2r-1};\mathbf z) 
 =Q^r(\mathbf z) P_1(\mathbf z)
 = S^{2r+1}(\mathbf z).
\end{equation}
 If $m=2r$ (for $r \geq 1$), then \eqref{normalaveraging} gives 
\begin{equation}%keepitlabeled 
 \lim_{N\to\infty}\frac{1}{N}
 \sum_{i=1}^{N}
 M(x_{2i}\cdots x_{2i+2r-1};\mathbf z) 
 = P_2(\mathbf z) Q^{r-1}(\mathbf z) P_1(\mathbf z)
 = S^{2r}(\mathbf z). 
\end{equation} 
 If $m=2r+1$ (for $r \geq 0$), then \eqref{normalaveraging} gives 
\begin{equation}\label{evenstartoddlength}
 \lim_{N\to\infty}\frac{1}{N}
 \sum_{i=1}^{N}
 M(x_{2i}\cdots x_{2i+2r};\mathbf z) 
 = P_2(\mathbf z) Q^r(\mathbf z) 
 = S^{2r+1}(\mathbf z).
\end{equation} 
 So, regardless of whether a length-$m$ block begins at an odd position or at an even position, the limit of each 
 displayed average among \eqref{oddstartevenlength}--\eqref{evenstartoddlength} 
 is $S^m(\mathbf z)$.  Moreover, among 
 the $n-m+1$ possible starting positions of length-$m$ blocks, the numbers of odd and 
 even starting positions are 
\begin{equation}\label{dividebyn} 
 \left\lceil\frac{n-m+1}{2}\right\rceil
 \quad\text{and}\quad
 \left\lfloor\frac{n-m+1}{2}\right\rfloor, 
\end{equation}
 respectively, and, furthermore, since $m$ is fixed, each of these quantities in \eqref{dividebyn}, after division by $n$,
tends to $\frac{1}{2}$. It follows from this and from 
 \eqref{oddstartevenlength}--\eqref{evenstartoddlength} that 
\begin{equation}\label{allstartpositions}
 \lim_{n\to\infty}\frac{1}{n}
 \sum_{j=1}^{n-m+1}M(x_jx_{j+1}\cdots x_{j+m-1};\mathbf z)
 = S^m(\mathbf z) 
\end{equation}
 for every fixed $m\geq1$. 

 Now, let $E$ be a (nonempty) decimal block of length $m$.   By the definition of $B_E$, the coefficient of     $z_0^{\nu_0(E)}  
  z_1^{\nu_1(E)}\cdots z_9^{\nu_9(E)}$ in the polynomial   $$  \frac{1}{n}
 \sum_{j=1}^{n-m+1}M(x_jx_{j+1}\cdots x_{j+m-1};\mathbf z)
$$ is $\frac{B_E(\Xi_{10},n)}{n}$. On the other hand,
\begin{equation}\label{takecoefficient}
 \left[z_0^{\nu_0(E)}z_1^{\nu_1(E)}\cdots z_9^{\nu_9(E)}\right]
 S(\mathbf z)^m
 =\frac{m!}{\nu_0(E)!\nu_1(E)!\cdots\nu_9(E)!}\frac{1}{10^m}.
\end{equation}
Since
\begin{equation}\label{equivalenceclasssize}
 \left|[E]_{\sim}\right|
 =\frac{m!}{\nu_0(E)!\nu_1(E)!\cdots\nu_9(E)!},
\end{equation}
taking the  coefficient indicated   in \eqref{allstartpositions} gives us that  
$$
 \lim_{n\to\infty}
 \frac{1}{\left|[E]_{\sim}\right|}
 \frac{B_E(\Xi_{10},n)}{n}
 =\frac{1}{10^m}
 =\frac{1}{10^{\ell(E)}}.
$$
   We thus have that   $\Xi_{10}$ is abelian-normal in base $10$.  

 So, it remains to show that $\Xi_{10}$ is not normal. In this direction, we begin by considering the ordered decimal block $01$. At starting 
 positions of the form $2i-1$, the block of length $2$ is exactly $\phi(c_i)$. The word $01$ has two
preimages under $\phi$, namely the base-$100$ symbols $1$ and $10$.
Since every base-$100$ symbol has limiting frequency $1/100$, we have
\begin{equation}\label{aligned01frequency}
 \lim_{N\to\infty}\frac{1}{N}
 \#\{1\leq i\leq N:x_{2i-1}x_{2i}=01\}
 =\frac{2}{100}.
\end{equation}
 At starting positions of the form $2i$, the relevant block is $\phi_2(c_i)\phi_1(c_{i+1})$. By \eqref{normalaveraging} with $s=2$, and by 
 \eqref{singlecoordinateidentities}, its limiting frequency is
\begin{align}
\begin{split}
 &\frac{1}{100^2}
 \#\{(u,v)\in\{0,1,\ldots,99\}^2:
 \phi_2(u)=0,\ \phi_1(v)=1\} \\
 &\qquad=
 \frac{1}{100^2}
 \#\{u:\phi_2(u)=0\}
 \#\{v:\phi_1(v)=1\}
 =\frac{10\cdot10}{100^2}
 =\frac{1}{100}.
\end{split}\label{crossing01frequency}
\end{align}
  Finally, starting positions of the forms $2i-1$ and $2i$ each have  asymptotic density $\frac{1}{2}$.  Combining
 \eqref{aligned01frequency} and \eqref{crossing01frequency}, we obtain
$$
 \lim_{n\to\infty}\frac{A_{01}(\Xi_{10},n)}{n}
 =\frac{1}{2}\frac{2}{100}
 +\frac{1}{2}\frac{1}{100}
 =\frac{3}{200}, 
$$
 in contrast to the value $\frac{1}{100}$   required for the frequency for the
block $01$ in a normal base-$10$ expansion. We thus have that $\Xi_{10}$ is not
normal in base $10$.
\end{proof}

\end{document}
